% !TEX TS-program = xelatex
\documentclass[11pt,a4paper]{article}

% ---------- Page ----------
\usepackage[a4paper,margin=1.8cm]{geometry}

% ---------- Fonts ----------
\usepackage{fontspec}
\setmainfont{Calibri}

% ---------- Colors ----------
\usepackage{xcolor}
\definecolor{PrimaryBlue}{RGB}{0,70,140}
\definecolor{SecondaryBlue}{RGB}{90,120,160}
\definecolor{RuleGray}{RGB}{210,210,210}

% ---------- Links ----------
\usepackage{hyperref}
\hypersetup{
  colorlinks=true,
  urlcolor=PrimaryBlue,
  linkcolor=PrimaryBlue
}
\urlstyle{same}

% ---------- Icons ----------
\usepackage{fontawesome5}
\newcommand{\hicon}[2][]{%
  \textcolor{PrimaryBlue}{\raisebox{-0.15ex}{\mdseries\faIcon[#1]{#2}}}
}

% ---------- Paragraph spacing ----------
\setlength{\parindent}{0pt}
\setlength{\parskip}{8pt}

% ---------- Section title style (same as resume) ----------
\newcommand{\cvsection}[1]{%
  \vspace{8pt}
  {\Large\bfseries\color{PrimaryBlue}#1}\hspace{10pt}%
  {\color{PrimaryBlue}\leaders\hrule height 1.2pt\hfill}\par
  \vspace{10pt}
}

\begin{document}

% ===================== Header =====================
{\fontsize{24pt}{26pt}\selectfont\bfseries\color{PrimaryBlue} Novin}\\[4pt]
{\large Mechanical Engineer — CFD \& FEA and Project Management}\\[6pt]

\small
\hicon{map-marker-alt}~\href{https://maps.google.com/?q=1+Rue+de+Trèves,+L-2631+Luxembourg}{1, Rue de Trèves, L-2631 Luxembourg} \quad
\hicon{envelope}~\href{mailto:mnovinnam1997@gmail.com}{mnovinnam1997@gmail.com}\\[3pt]
\hicon{phone}~\href{tel:+352661901293}{+352 661 901 293} \quad
\hicon{linkedin-in}~\href{https://linkedin.com/in/novinnam}{linkedin.com/in/novinnam}
\normalsize

\vspace{14pt}

\today

\vspace{10pt}

Rotarex S.A.\\
Luxembourg\\

\vspace{14pt}

\cvsection{Cover Letter}

Dear Rotarex S.A. Team,

I saw your message in your website about AI generating tools, OK here is my cover letter in
my own words.

My background is fully technical in mechanical engineering, from bachelor which I got familiar
with CFD to master with FEA focus.

I have done my master thesis in PAUL WURTH company in Luxembourg. The project was implementing
a camera enclosure nozzle integrated with air knife to protect the camera from dusty environment.

As a master student this work started from research in academic papers to existed products in
the market.

I designed the prototypes with CREO, however I have learnt computer aided design with
different tools liks CATIA, Solidworks in bachelor and Fusion 360 and Invetor in master, since
the features inside these applications are the same, I could easily adapt to CREO.

After studying different approaches, and dust condition in industrial environment in the site,
I applied k-omega SST to the air side and lagrangian particle tracking to the dust side in ANSYS.

So the drag coefficients and equations were chosen based based on physics of my problem.
The real challenge was random behavior of Bambu lab P1S 3D printer, after asking and searching,
I could not find the official technical specification for tolerance in this printer, so that based
on experiences I tried to get the gap size I wanted to be close to my calculation which was 0.05mm,
 however, I end with 3 different prototypes to make a comparison.
 
The validation section has been done in lab with 2 main parts: first part was quantitative by
checking velocity of air in different probe positions. The second was qualitative by checking
the behavior of dust particles and smoke machine in different angle of attack.

As a result from simulation to manufacturability is covered in my master thesis.

About my master I got familiar with FEA, specifically in machine design exercise course, I have
applied VON MISES into single power screw press machine and with topology optimization I tried
to reduce the mass in that machine.

I am not talking more, my resume is attached for your reference.
Simply, I am applying to this position because I am junior level mechanical engineer and 
nothing else I can mention how I am motivated and related into this position.


\vspace{12pt}

Yours sincerely,\\[10pt]
Novin

\end{document}